% ============================================================
% 第 9 章  电力系统的负荷
% （教材：何仰赞、温增银《电力系统分析》下册 第 9 章）
% 结构对齐 9.1—9.3；例题、习题不录
% ============================================================

\chapter{电力系统的负荷}

前面八章把发电机、变压器、线路都变成了等值电路，剩下的最后一个元件是负荷。
负荷是系统的“目的”，也是分析中最难精确描述的部分：它由成千上万台设备组成，
随生产、作息、气候而变，且对电压、频率都有依赖。
本章给出三个层次的内容：\textbf{负荷是什么}（组成）、\textbf{负荷怎样随时间变}（曲线与指标）、
\textbf{负荷怎样随电压频率变}（特性与模型）。

\section{负荷的组成}

系统中所有电力用户的用电设备所消耗的电功率总和，就是电力系统的负荷，
亦称\textbf{综合用电负荷}——它把不同地区、不同性质的所有用户负荷总加起来得到。

\subsection{用电设备的类型与用户分类}

主要用电设备有：异步电动机、同步电动机、电热装置、整流装置、照明设备。
按用户性质可分为工业负荷、农业负荷、交通运输业负荷和人民生活用电负荷。
在不同性质的用户中，各类设备消耗功率的比重不同；在工业内部，不同行业的比重也不同。

\begin{table}[htbp]
  \centering
  \caption{几个工业部门用电设备比重的统计（单位：\%）}
  \small
  \begin{tabular}{@{}lcccccc@{}}
    \toprule
    用电设备 & 综合性中小工业 & 纺织工业 & 化学工业（化肥、焦化） & 化学工业（电化） & 大型机械加工 & 钢铁工业 \\
    \midrule
    异步电动机 & 79.1 & 99.8 & 56.0 & 13.0 & 82.5 & 20.0 \\
    同步电动机 & 3.2 & — & 44.0 & — & 1.3 & 10.0 \\
    电热装置 & 17.7 & 0.2 & — & — & 15.0 & 70.0 \\
    整流装置 & — & — & — & 87.0 & 1.2 & — \\
    \midrule
    合计 & 100.0 & 100.0 & 100.0 & 100.0 & 100.0 & 100.0 \\
    \bottomrule
  \end{tabular}
\end{table}

这张表说明一件重要的事：\textbf{钢铁、电化等行业的负荷以电热、整流为主，功率因数高、
对电压不敏感；纺织、机械行业以异步电动机为主，功率因数低、无功占比大}。
负荷构成决定了无功需求，也决定了第 12 章调压措施的取舍。

\subsection{负荷的功率因数与无功构成}

异步电动机在总负荷中占比大（多数行业 $50\%\sim90\%$），其功率因数随负载率变化剧烈：
\begin{itemize}
  \item 额定负载时 $\cos\varphi\approx0.85\sim0.90$；
  \item 半载时 $\cos\varphi\approx0.7\sim0.8$；
  \item 空载时 $\cos\varphi$ 只有 $0.2\sim0.3$（几乎全是励磁无功）。
\end{itemize}
因此“大马拉小车”（电动机长期轻载）是系统无功需求偏大、功率因数偏低的重要原因。
这也是无功补偿（第 12 章）与“电动机经济运行”的直接依据。

系统的自然功率因数一般在 $0.7\sim0.85$。若不做补偿，无功负荷将大量占用发电机容量与
输电能力：输送同样的有功，功率因数从 $0.85$ 降到 $0.7$，电流增大约 $21\%$，
线路损耗增大约 $47\%$。这就是“无功就地平衡”原则的经济含义。

\subsection{三级负荷概念}

\begin{table}[htbp]
  \centering
  \caption{综合用电负荷、供电负荷与发电负荷}
  \small
  \begin{tabular}{@{}p{0.20\linewidth}p{0.34\linewidth}p{0.34\linewidth}@{}}
    \toprule
    名称 & 定义 & 关系 \\
    \midrule
    综合用电负荷 & 所有用户用电设备消耗的功率之和 & 基准量 \\
    供电负荷 & 综合用电负荷 $+$ 电力网的功率损耗 & 各发电厂应供给的功率 \\
    发电负荷 & 供电负荷 $+$ 发电厂厂用电 & 各发电厂应发出的功率 \\
    \bottomrule
  \end{tabular}
\end{table}

三者的差额就是“网损”与“厂用电”，它们正是第 14 章经济运行的考核对象
（网损率、厂用电率）。

\section{负荷曲线}

负荷随时间变化的规律用\textbf{负荷曲线}描述。常用的有日负荷曲线和年负荷曲线。

\subsection{日负荷曲线}

日负荷曲线描述一天 24 小时负荷的变化。曲线最大值称\textbf{日最大负荷} $P_{\max}$（峰荷），
最小值称\textbf{日最小负荷} $P_{\min}$（谷荷）。为便于计算，常把连续曲线画成阶梯形。

由日负荷曲线可算一日总耗电量与日平均负荷：
\keyformula{W_d=\int_{0}^{24}P\,\mathrm{d}t,\qquad P_{av}=\frac{W_d}{24}}
衡量曲线起伏特性的两个系数：
\begin{align}
  \text{负荷率}\quad k_m&=\frac{P_{av}}{P_{\max}},\\
  \text{最小负荷系数}\quad \alpha&=\frac{P_{\min}}{P_{\max}}.
\end{align}
峰谷差 $\Delta P_m=P_{\max}-P_{\min}$ 是调峰容量的依据。
这两个系数不仅用于日负荷曲线，也用于其他时段的负荷曲线。

不同性质用户的日负荷曲线形状差别很大：
\begin{itemize}
  \item 三班制连续生产的重工业（如钢铁）：曲线平坦，$\alpha$ 可达 $0.85$；
  \item 一班制生产的轻工业（如食品）：变化幅度大，$\alpha$ 只有 $0.13$ 左右；
  \item 非排灌季节的农业：农村加工用电每天仅约 12 小时；
  \item 市政生活：有明显照明高峰。
\end{itemize}
由于各用户的最大、最小负荷不在同一时刻出现，
\textbf{系统的最大负荷总小于各用户最大负荷之和，系统的最小负荷总大于各用户最小负荷之和}
——这就是“负荷同时率”现象的来源，也是系统调峰比单用户调峰更容易的原因。

日负荷曲线是安排\textbf{日发电计划}和确定\textbf{系统运行方式}的重要依据。

\subsection{峰谷差与调峰}

日负荷曲线最受关注的特征是\textbf{峰谷差}
\begin{equation}
  \Delta P_m=P_{\max}-P_{\min}=P_{\max}\left(1-\alpha\right).
\end{equation}
峰谷差决定了系统需要多少“可调节容量”：峰荷时要有足够的机组顶峰，
谷荷时要有足够的机组压负荷（或停机）。若调节能力不足，只能靠“开停机”应对，
既不经济也影响设备寿命。

\begin{table}[htbp]
  \centering
  \caption{典型用户的负荷曲线特征}
  \small
  \begin{tabular}{@{}p{0.24\linewidth}p{0.18\linewidth}p{0.46\linewidth}@{}}
    \toprule
    用户类型 & 最小负荷系数 $\alpha$ & 曲线特征 \\
    \midrule
    三班制重工业（钢铁） & 约 $0.85$ & 平坦，基荷型，宜带基荷 \\
    一班制轻工业（食品） & 约 $0.13$ & 峰谷差大，昼开夜停 \\
    农业（非排灌季节） & — & 每天约 12 小时用电，季节性明显 \\
    市政生活 & — & 有明显照明晚高峰 \\
    \bottomrule
  \end{tabular}
\end{table}

系统峰谷差与用户构成有关：工业比重大则曲线平坦，生活用电比重大则峰谷差大。
随着新能源（光伏）大规模接入，\textbf{“鸭子曲线”}现象（午间光伏出力大、
净负荷下凹，傍晚光伏退出、负荷上升）使峰谷差进一步扩大，对调峰能力提出更高要求。
这也是第 13 章“调频与调峰”、第 14 章“经济运行”要解决的问题的现实背景。

\subsection{年最大负荷曲线与年持续负荷曲线}

\paragraph{年最大负荷曲线}
描述一年内每月（或每日）最大有功功率负荷的变化，主要用来安排\textbf{发电设备的检修计划}，
也为制订发电机组或发电厂的扩建、新建计划提供依据。

\paragraph{年持续负荷曲线}
按一年中系统负荷的数值大小及其持续小时数顺序排列绘制。例如全年 8760 小时中，
有 $t_1$ 小时负荷为 $P_1$（$=P_{\max}$），$t_2$ 小时为 $P_2$，……，于是得到阶梯下降的曲线。
由它可以算出全年耗电量
\begin{equation}
  W=\int_{0}^{8760}P\,\mathrm{d}t .
\end{equation}
若负荷始终等于 $P_{\max}$，经过 $T_{\max}$ 小时所消耗的电能恰好等于全年实际耗电量，则
\keyformula{T_{\max}=\frac{W}{P_{\max}}=\frac{1}{P_{\max}}\int_{0}^{8760}P\,\mathrm{d}t}
$T_{\max}$ 称\textbf{最大负荷利用小时数}。它是衡量负荷特性的综合性指标：
$T_{\max}$ 越大，说明负荷越平稳、设备利用率越高。年持续负荷曲线常用于安排发电计划
和进行可靠性估算。

\begin{table}[htbp]
  \centering
  \caption{负荷曲线的种类与用途}
  \small
  \begin{tabular}{@{}p{0.22\linewidth}p{0.24\linewidth}p{0.42\linewidth}@{}}
    \toprule
    曲线 & 描述对象 & 主要用途 \\
    \midrule
    日负荷曲线 & 一天 24 小时 & 安排日发电计划、确定运行方式、调峰 \\
    年最大负荷曲线 & 每月（日）最大值 & 安排检修计划、规划扩建 \\
    年持续负荷曲线 & 负荷大小与持续小时数 & 年耗电量、$T_{\max}$、可靠性估算 \\
    \bottomrule
  \end{tabular}
\end{table}

\corepoint{三个指标不要混：\textbf{负荷率} $k_m=P_{av}/P_{\max}$ 反映“平均用得多满”；
\textbf{最小负荷系数} $\alpha=P_{\min}/P_{\max}$ 反映“低谷有多低”（调峰难度）；
\textbf{最大负荷利用小时数} $T_{\max}=W/P_{\max}$ 反映“全年等效满负荷运行多少小时”。
$k_m$ 与 $T_{\max}$ 的关系是 $T_{\max}=8760\,k_m$。}

\section{负荷特性与负荷模型}

\subsection{负荷的静态特性}

负荷特性指负荷功率随电压、频率变化而变化的规律。
\textbf{静态特性}指电压、频率缓慢变化时的规律；\textbf{动态特性}指快速变化（如短路、电动机
失速）时的规律。本课程主要用静态特性。

\paragraph{电压静态特性（二次多项式，ZIP 模型）}
\begin{equation}
  P=P_N\left[a_p\left(\frac{U}{U_N}\right)^{2}+b_p\left(\frac{U}{U_N}\right)+c_p\right],
\end{equation}
\begin{equation}
  Q=Q_N\left[a_q\left(\frac{U}{U_N}\right)^{2}+b_q\left(\frac{U}{U_N}\right)+c_q\right],
\end{equation}
系数由实测特性用最小二乘法拟合，且满足
\begin{equation}
  a_p+b_p+c_p=1,\qquad a_q+b_q+c_q=1 .
\end{equation}
三项各有明确物理含义：
\begin{itemize}
  \item 与 $U^{2}$ 成正比的第一项：代表\textbf{恒定阻抗}消耗的功率；
  \item 与 $U$ 成正比的第二项：代表\textbf{恒定电流}负荷对应的功率；
  \item 常数项：代表\textbf{恒定功率}分量。
\end{itemize}
故该模型又称 \textbf{ZIP}（恒阻抗—恒电流—恒功率）模型。

\paragraph{线性化模型}
当电压、频率在额定值附近小幅度变化时，可线性化：
\begin{align}
  P&=P_N\left(1+k_{pv}\delta U\right),&
  Q&=Q_N\left(1+k_{qv}\delta U\right),\\
  P&=P_N\left(1+k_{pf}\delta f\right),&
  Q&=Q_N\left(1+k_{qf}\delta f\right),
\end{align}
其中
\begin{equation}
  \delta U=\frac{U-U_N}{U_N},\qquad \delta f=\frac{f-f_N}{f_N}.
\end{equation}
$k_{pv}$、$k_{qv}$ 称\textbf{负荷的电压调节效应系数}，$k_{pf}$、$k_{qf}$ 称\textbf{频率调节效应系数}。
同时考虑电压与频率变化时可用
\begin{align}
  P&=P_N\left(1+k_{pv}\delta U\right)\left(1+k_{pf}\delta f\right),\\
  Q&=Q_N\left(1+k_{qv}\delta U\right)\left(1+k_{qf}\delta f\right).
\end{align}
这些系数是第 12 章（电压与无功）和第 13 章（频率与有功）中“负荷自身调节作用”的量化依据。

\subsection{负荷的动态模型}

反映动态特性的模型一般由微分方程与代数方程组成，随所研究问题不同而取不同形式。
建立综合负荷动态模型包含两件事：\textbf{制订模型结构}与\textbf{确定模型参数}。
由于综合负荷所代表的设备数量极大、分布极广、种类繁多，工作状态带有随机性与时变性
（甚至跃变性），配电网结构也可能变化，因此“既准确又实用”的负荷模型至今仍未很好解决。

\subsection{综合负荷的等值电路}

电力系统分析中元件都用等值电路代表，负荷也不例外。最常采用的综合负荷等值电路有：
\begin{itemize}
  \item \textbf{含源等值阻抗（或导纳）支路}；
  \item \textbf{恒定阻抗（或导纳）支路}；
  \item \textbf{异步电动机等值电路}（阻抗值随转差而变）；
  \item 上述电路的不同组合。
\end{itemize}

\begin{table}[htbp]
  \centering
  \caption{不同计算中负荷模型的取用}
  \small
  \begin{tabular}{@{}p{0.18\linewidth}p{0.34\linewidth}p{0.38\linewidth}@{}}
    \toprule
    计算类型 & 负荷的处理 & 理由 \\
    \midrule
    潮流计算 & 恒定功率 $S_{LD}=P_{LD}+jQ_{LD}$（必要时用线性化静态特性） & 负荷功率由用户决定，与网络电压弱相关 \\
    短路计算 & 含源阻抗支路，或恒定阻抗支路 & 短路瞬间电动机可倒送电流（第 6 章） \\
    稳定计算 & 恒定阻抗，或恒定阻抗与异步电动机的组合 & 需反映电压、频率变化时的功率响应 \\
    \bottomrule
  \end{tabular}
\end{table}

采用“恒定阻抗 $+$ 异步电动机”组合模型时，还须补充一个反映异步电动机机械运动状态的
\textbf{转子运动方程}（第 17 章），否则无法描述电动机失速与恢复的过程。

\corepoint{同一套系统，在潮流、短路、稳定三类计算中负荷模型不同，原因不是“简化程度不同”，
而是\textbf{所关心的物理过程不同}：潮流关心功率分布（负荷按功率给定）、
短路关心瞬间电流（负荷可视为阻抗或电源）、稳定关心机电过程（负荷须有动态）。
选模型前先问“我要算的是哪个时间尺度上的什么量”。}
